\begin{exsol}{23}{Tadelis Exercise 11.5}{236--237}{Exam}{L4}
\asked Three-period alternating-offer bargaining over a pie of 1, player~1 proposing first, with
discount factors $\delta_1 \ne \delta_2$: (a) the SPE outcome; (b) player~1's advantage when
$\delta_1 = \delta_2$; (c) when player~2 has the advantage, and why.

\method Backward induction in finite alternating-offer bargaining (Tadelis \S11.2, pp.~224--227;
Gibbons \S2.1.D, pp.~68--70): in the last period the proposer takes everything; in each earlier
period the proposer offers the responder exactly the discounted value of what the responder would
get next period, and the responder accepts at indifference (Prop.~11.2's argument, p.~223).

\solution
\textbf{Notation.} $x_t$ = player~1's share of the pie when agreement is reached in period
$t \in \{1,2,3\}$; player $i$'s payoff from agreement $(x, 1-x)$ in period $t$ is $\delta_i^{t-1}$
times his share. Players~1, 2, 1 propose in periods 1, 2, 3; no agreement in period~3 gives $(0,0)$.

\textbf{Step 1: period 3.} Player~1 offers $x_3 = 1$ and player~2 accepts (she gets 0 either way).
Discounted payoffs: $(\delta_1^2,\ 0)$.

\textbf{Step 2: period 2.} Player~2 must give player~1 at least what he gets by rejecting:
\begin{equation}
  \delta_1\, x_2 \;\ge\; \delta_1^2\, x_3 = \delta_1^2 \;\Rightarrow\; x_2 = \delta_1 . \label{eq:23x2}
\end{equation}
Player~2 keeps $1 - x_2 = 1 - \delta_1$, worth $\delta_2(1 - \delta_1)$ in period-1 terms. (She prefers
this to making an offer player~1 rejects, which gives her 0.)

\textbf{Step 3: period 1.} Player~1 must give player~2 at least her period-2 value:
\[
  1 - x_1 \;\ge\; \delta_2\,(1 - x_2) = \delta_2\,(1 - (\delta_1)) ,
\]
so, keeping as much as possible,
\begin{equation}
  x_1 = 1 - \delta_2(1 - \delta_1) = 1 - \delta_2 + \delta_1\delta_2 . \label{eq:23x1}
\end{equation}
\textbf{(a)} The SPE outcome is immediate agreement on
\[
  (v_1, v_2) = \bigl(1 - \delta_2(1 - \delta_1),\ \ \delta_2(1 - \delta_1)\bigr).
\]
(Strategies: player~1 offers $x_1$ in period~1 and $1$ in period~3, and accepts in period~2 iff
offered at least $\delta_1$; player~2 accepts in period~1 iff offered at least
$\delta_2(1-\delta_1)$, offers $x_2 = \delta_1$ in period~2, accepts anything in period~3.)

\textbf{(b) Equal discount factors.} Setting $\delta_1 = \delta_2 = \delta$ in \eqref{eq:23x1}:
\[
  v_1 = 1 - \delta + \delta^2, \qquad v_2 = \delta - \delta^2,
\]
\[
  v_1 - v_2 = 1 - 2\delta + 2\delta^2 = (1 - \delta)^2 + \delta^2 \;>\; 0 .
\]
So player~1 gets strictly more than half for every $\delta$: he proposes first \emph{and} last.

\textbf{(c) When player 2 is ahead.} Player~2 gets more than half iff
\[
  \delta_2\,(1 - \delta_1) > \tfrac12 \iff \delta_1 < 1 - \frac{1}{2\delta_2} ,
\]
which needs $\delta_2 > \frac12$ (a patient player~2) and $\delta_1$ low enough (an impatient
player~1). Example: $\delta_1 = 0.2$, $\delta_2 = 0.9$ gives $v_2 = 0.9 \times 0.8 = 0.72$.
\emph{Why.} Player~1's power comes from the last period, but an impatient player~1 values it at only
$\delta_1$ from period~2's viewpoint, so player~2 keeps $1 - \delta_1$ when she proposes. A patient
player~2 values that period-2 position almost fully ($\delta_2$), so in period~1 player~1 must
concede $\delta_2(1 - \delta_1)$. Patience is bargaining power: $v_1$ rises with $\delta_1$
($\partial v_1/\partial\delta_1 = \delta_2 > 0$) and falls with $\delta_2$
($\partial v_1/\partial \delta_2 = -(1 - \delta_1) < 0$), as in the infinite-horizon formula
$\frac{1 - \delta_2}{1 - \delta_1\delta_2}$ (Tadelis \S11.3; \texttt{Notas/L04/L04-07}).

\interp First-mover and last-mover advantages favour player~1, so only a large patience gap in
player~2's favour can overturn them. With equal patience player~1's edge $(1-\delta)^2 + \delta^2$
is smallest at $\delta = \frac12$ (edge $\frac12$, $v_1 = \frac34$). As $\delta \to 0$ the first-mover
advantage takes over and $v_1 \to 1$; as $\delta \to 1$ the last-mover advantage takes over and again
$v_1 \to 1$, because waiting for period~3 costs player~1 nothing.

\verify \texttt{check\_B\_L04.py} redoes the recursion symbolically ($x_1 = 1 - \delta_2 +
\delta_1\delta_2$; $v_1 - v_2 = (1-\delta)^2 + \delta^2$), tests the condition in (c) on four
$(\delta_1, \delta_2)$ pairs, and repeats the backward induction numerically with offers on a grid of
step $\frac1{1000}$ for three pairs, e.g.\ $(\delta_1,\delta_2) = (0.9, 0.7)$ gives $x_1 = 0.93$, the
value in \texttt{Notas/L04/L04-10-formulas.md}. Tadelis prints no answer.

\answer{(a) Immediate agreement: player~1 gets $1 - \delta_2(1 - \delta_1)$, player~2 gets
$\delta_2(1 - \delta_1)$. (b) With $\delta_1 = \delta_2 = \delta$: $v_1 - v_2 = (1-\delta)^2 + \delta^2
> 0$. (c) Player~2 is ahead iff $\delta_2(1 - \delta_1) > \frac12$: player~2 patient
($\delta_2 > \frac12$) and player~1 impatient ($\delta_1 < 1 - \frac1{2\delta_2}$), because player~1's
last-period power is worth little to an impatient player~1, while player~2's period-2 power is worth a
lot to a patient player~2.}
\end{exsol}
